\section{Christendom and the Superstition of Fact}

It seems necessary at this time to make some ideas clear, as least, clearer. First of all, metaphysics is not the same thing as philosophy. Metaphysics is knowing, an activity of spirit, active and masculine. Philosophy is thinking, an activity of the soul, passive, and feminine. A man knows himself to be free, with an active will; he has no need to wade through several hundred pages of Kant to convince himself of what he knows in his own immediacy.

A warrior acts courageously because it is an expression of his nature, his being. An artist can portray a courageous man in a poem or play, but is not thereby courageous himself. A philosopher will debate the meaning of courage, but will always end in aporia. If a scout comes back to camp with news that the enemy is on the march, a debate about the meaning of the other will not then ensue. Nor will anyone attempt to achieve the synthesis between the self and the other as enemy.

When we write of doctrines, we mean metaphysical doctrines, not specifically religious or theological doctrines. Here are some quotes from \emph{Spiritual Authority and Temporal Power} by Rene Guenon that clarify Gornahoor's point of view:

\begin{quotex}
During the Middle Ages there existed throughout the West a real unity, based on properly traditional foundations, which we call `Christendom', but when these secondary unities of a purely political --- that is to say temporal and no longer spiritual --- order were formed this great unity of the West was irremediably broken and the effective existence of Christendom came to an end. Nations, merely the dispersed fragments of what was formerly Christendom, false unities substituted for the true one by the temporal power's will to dominate can, given the very conditions of their origin, survive only by opposing each other and ceaselessly contending among themselves in all fields. Now \textbf{spirit is unity}, \textbf{matter is multiplicity and division}; and the more one removes oneself from spirituality, the more antagonisms are accentuated and amplified.
\end{quotex}


The last sentence speaks for itself. Please understand the following point thoroughly.

\begin{quotex}
When speaking of Catholicism the utmost care must always be taken to distinguish between what conerns Catholicism itself as a doctrine and what relates only to the present organizational state of the Catholic church. Whatever one may think about the second, it cannot affect the first. ... there are very few today who are able, when such a need arises, to free themselves from historical contingencies, to the extent that certain defenders of Catholicism, not only its adversaries, believe that everything can be reduced to a simple question of historicity, which is one form of the modern \textbf{superstition of fact}.
\end{quotex}


As for the continuity of the Roman Tradition, Guenon points out

\begin{quotex}
The symbolic assimilation of Christ with Janus as the supreme prinnciple of the two powers is the very clear sign of a \textbf{certain traditional continuity} (\emph{too often ignored or deliverately denied}) \textbf{between ancient Rome and Christian Rome}; and we must not forget that in the Middle Ages the empire was just as Roman as the papacy.
\end{quotex}


\flrightit{Posted on 2011-01-27 by Cologero}

\begin{center}* * *\end{center}

\begin{footnotesize}
\begin{sffamily}
\texttt{Perennial on 2011-01-28 at 02:34 said:}

Great article! The material here is getting better and better!

\hfill

\texttt{Mark on 2011-01-28 at 02:40 said:}

I will start out by stating that I accept the view of Guenon in \textit{Introduction to the Study of the Hindu Religions} about the character of Vedanta as a Pure Metaphysics that is elaborated in the chapter titled \textit{Essential Character of Metaphysics}, and I will quote part of the first sentence.

``While the religious point of view necessarily implies the intervention of an element drawn from the sentimental order, the metaphysical point of view is exclusively intellectual''.

So, let's not be concerned with sentimental attachments to paganism, Islam, Catholicism, etc.. Let us seek after what embodies Metaphysics. For this I will use Vedantic notions of the Brahman.

Definition: Brahman is Satyam (Truth), Jnanam (Knowledge), Anantam (Infinity), Anandam (Bliss), and Amalatava (Unstainable Purity). I have read this before as Satyam being perfect, unchangeable Being, Jnanam being Consciousness that is unrestricted, Anantam being unconditioned, not limited by space and time, for Anandam, bliss is a good enough description, and Amalatava means unstained by karma. So here we see a concept of pure metaphysics of the unconditioned, and it is a key doctrine of the Advaita school, which is the favorite of Guenon that the highest state is the merging of the Atman and Brahman.

Here we see a metaphysics of what I like to call ``Transtheism'', which is the proper position of the Indo-European religions. The gods are not denied, but what transcends them is a state of Absolute Being, which we can attain. Here is my argument.

1) The ``Transtheism'' that is stated by Vedanta is pure Metaphysics

2) The Indo-European systems that we would call colloquially as Slavic paganism, Norse paganism, Baltic paganism, Hellenic paganism, and the other Indo-European paganisms are properly transtheistic

3) Any form of Christianity cannot be transtheistic, given its Abrahamic character it is monotheistic. God will always be seen as distinct from man, and the contingent and dependent nature of man upon a distinct being (God) cannot be transcended.

4) Accepting the notion of what is syncretic vs synthesis as to deal with what is congruent vs incongruent, then the ``paganisms'' synthesize with the Vedantic metaphysics, since they are congruent, whereas Christianity is syncretic with the Vedantic metaphysics, since it is incongruent with it.

5) Therefore systems like Rodnovery and Asatru are a superior choice to Christianity since they synthesize with Vedanta, as opposed to being syncretic like Christianity, since it is incongruent with the Metaphysics of Vedanta.

6) The only choice for Christianity over one of our pre-Christian traditions is due to sentimentalism, and not Metaphysics.

\hfill

\texttt{Perennial on 2011-01-28 at 02:52 said:}

I am not sure about that Mark. In Eastern theology we often see theosis emphasis, of man becoming a god, so to speak. It has been observed that Christ came and died so that we may be like Him by realizing that oneness, and union, with Him who is Eternal, that is God. In other words, Christ came to show how we may be Him. Shedding the ignorance of our material existence, while not condemning it, will help us to realize our transcendent nature and come into gnosis of Eternal Truth. The Sacraments are just an exterior means of finding this Truth (i.e. ``a means of grace''), just like meditation and ascetic practice. As means, they are a help. When they become ends in themselves, then the inner signifigance is lost, and the Sacrament becomes merely exoteric, as it has today. The Tridentine Mass was lost because the entire foundation of why it was built was lost to a utilitarian, populist, and democratic theology. But turning God into an object of worship does not make Him an idol.

\hfill

\texttt{Mark on 2011-01-28 at 03:05 said:}

So, do you deny 3? Tell me the propositions that you deny, and why?

\hfill

\texttt{Perennial on 2011-01-28 at 03:37 said:}

It is somewhat difficult to find the right words to describe it. This is why Eastern theologians always emphasized being apophatic, or only addressing what God and theology is not, rather then is. But I will give my best:

Religion is often described to mean ``to bind back.'' All creation is seeking to come back to the One, the Holy Trinity. Creation is distinct because of the Fall, but duality did not exist for Adam \& Eve. In Him ``they lived, and moved, and had their being.'' The Fall brought ignorance, and thus duality, and loss of understanding. Thus, the narrative of man since the Fall is a ``binding back,'' since man must now rediscover that union that was lost. This is why Adam realized he was naked, because the higher Self was subsumed by the lower self, and in partaking of the forbidden fruit he actually rejected transendance in favour of the material. He did not recognize the material as ``seperate'' prior to the Fall. The knowledge of good and evil meant a loss of understanding. This is Original Sin. The loss of original unity, and the loss of transendence, is now the common inheritance of all men, and thus he cannot be ``saved'' without rising to the higher self, and thus by Holy Baptism he rejects ignorance and embraces the path toward his higher nature. As man thus progresses, in Soul, Spirit, and Sacrament, he loses the lower self and understands the Mysteries of the Divine, if he seeks them meetly. Eventually, the old self will die away, and he will Resurrect on the last day, the whole man, and be reconciled back unto the One, entering eternal and undivided union, in full knowledge of the Truth. He will not be absorbed into God, but will ``... clothe (y)ourselves with the Lord Jesus Christ, and not think about how to gratify the desires of the flesh.'' (St. Paul to the Roman 13:14) Christ thus came and died that we might follow Him, not as slavish worshipers, but ``If any man will come after Me, let him deny himself, and take up his cross, and follow me.'' (St. Matthew 16:24) To follow Christ is to be Christ, and do as He did, and sacrifice to lower self to the higher. ``For whosoever will save his life shall lose it; but whosoever shall lose his life for my sake and the Gospel's, the same shall save it.'' (St. Mark 8:33-35). This is how Christ is Our Saviour. He came to show the Way of Truth and Liberation, and this is to free us from material bondage and embark on the eternal. Only through gnosis can a man find the One, but through exoteric practice one at least begins the journey.

\hfill

\texttt{Perennial on 2011-01-28 at 03:41 said:}

So yes, I would say 3 is an error, since it is presumed, not proven, and in fact I think evidence shows otherwise. 4,5, \& 6 are therefore ipso facto erroneous, as they are based on 3. That being said, in order for this and 1 and 2 to be judged, I request that you expand upon your definition of ``Transtheism.'' I myself am a Panentheist (all things exist in God, and are therefore not outside of, or foreign to, Him).

\hfill

\texttt{Mark on 2011-01-28 at 04:27 said:}

I am not sure that you understand the metaphysics of Advaita Vedanta, so the Sanskrit terms that I use might not be understood, and this might be the problem why you do not understand what I mean by ``Transtheism''. The idea is that the Metaphysical ``God'' is the conditions that transcend existence. So, you are no longer conditioned in any way, you are no longer in the flux of becoming, but you are Satyam, you are no longer limited in intellect, but Jnanam, you are no longer conditioned by time, space, and causation, but Anantam, you are no longer tainted by Karma, but Amalatava, and you experience perfect bliss (Anandam).

In Christianity, you always remain a contingent being under the conditioning factor of a separate and distinct being (God). No Abrahamic religion can be transtheistic, since their deity is always a distinct being.

This is actually quite for wikipedia: \url{http://en.wikipedia.org/wiki/Transtheism}


Another reason that I like the term is that a rather good view of it is in Ride the Tiger in a chapter called Beyond Atheism and Theism

\hfill

\texttt{Perennial on 2011-01-28 at 06:35 said:}

I am not sure how what I described differs what you described. Perhaps you could elaborate? I speak of reaching liberation, where the fiction that seperates your from your Creator disappears. There is a Creator, and we are the created. Obviously, we are limited by the words we have, but essentially the Creator and the created share a like substance (God created us out of Himself, therefore we are ``in His Image'' and are like unto Him. We need to lose our duality and be with Him, in Him, in the Unity of the Holy Ghost. We are only seperate in nature (i.e. we are not God by design) but not in essence (the spirit returns to it's source, the One). When we become theosis, we will no longer be a seperate being, but ``Very God.'' We shall be a unique component of His Oneness, but shall be of it nonetheless.

So either I am not understanding what you are trying to say, or you not understanding me. We shall see.

\hfill

\texttt{Cologero on 2011-01-28 at 08:45 said:}

Please restrict your discussion to `Christendom' as Guenon defined it above, not some ill-defined ``Christianity''.

Gornahoor has addressed what is being called ``transtheism'' ... find it and discuss it please.

The typical pagan man in the street was a polytheist. The philosophical elite and their aristocratic followers were monotheists and did not even accept the existence of the gods.

Asatru is a fabricated system and quite modernist. The very idea of making a ``choice'' is modernist.

Boris Mouravieff wrote a monograph on Slavic paganism. Please locate it and discuss it if you like.

You can debate systems of belief all you want, but Gornahoor is focusing on states of consciousness that can be attained. Without that, all discussions are futile and result from a fundamentally materialist mindset.

\hfill

\texttt{Cologero on 2011-01-28 at 10:45 said:}

\textbf{What Buddha Suggested}
\begin{verse}
Do not be led by hearsay\\
or by what is handed down by tradition\\
or by what people say,\\
or by what is stated by the authority of your traditional teachings.\\
Be not led by reasoning,\\
nor by inferring,\\
nor by argument as to method,\\
nor by delight in speculative opinions,\\
nor by seeming possibilities,\\
nor by the directions from your teachers.

But, when you know of yourself\\
that certain actions done by you are not good,\\
wrong and considered worthless by the wise;\\
when followed and put in to practice,\\
lead to loss or suffering, then give them up ... \\
and when you know of yourself that\\
certain actions done by you are good, true\\
and considered worthy by the wise,\\
then accept them and put them into practice.

... ..

I have shown you the methods that lead to liberation.\\
But you should know that liberation depends upon yourself.
\end{verse}

\hfill

\end{sffamily}
\end{footnotesize}

